\chapter{超空間と超対称力学の構成}

\lectureoverview{超対称性が作用しうる最小の系、すなわち等しいエネルギーをもつ1つのボソン振動子と1つのフェルミオン振動子から超対称性を再構築する。両者の2階方程式と1階方程式の対比から超電荷とその代数が導かれる。次にその代数を、通常の座標(t)とグラスマン座標\(\theta,\bar\theta\)をもつ空間上の幾何学として再解釈する。超場はボソン、フェルミオン、補助場を単一の対象にまとめ、ベレジン積分は不変な通常のラグランジアンを抽出し、共変微分は対称性と整合する拘束を可能にし、カイラル拘束が最初の重要な例を与える。本講義はこの仕組みの符号と順序に関する困難を意図的に露呈させ、最後に、なぜ対称性そのものがファインマン則の目録よりも根本的であるかを説明して締めくくる。}

\section{一モード系に戻る}

静止したフェルミオンとボソンから再び始めよう。空間座標と運動量ラベルはすべて抑制する。量子は与えられたモードに存在するか存在しないかのいずれかであり、生成消滅演算子だけでそれを記述するのに十分である。ボソンモードを\(a^\dagger,a\)、フェルミオンモードを\(c^\dagger,c\)と書く。これらは
\begin{align}
[a,a^\dagger]&=1,
&
\{c,c^\dagger\}&=1,
&
c^2=(c^\dagger)^2&=0.
\end{align}
を満たす。2つの量子が同じエネルギー\(m\)をもつならば、無関係な真空定数を除いて
\begin{equation}
H=m\left(a^\dagger a+c^\dagger c\right).
\label{eq:l7-one-mode-hamiltonian}
\end{equation}
である。数演算子\(a^\dagger a\)は任意の非負整数値をとりうるが、\(c^\dagger c\)は0か1しかとれない。これがおもちゃ模型の運動学のすべてである。

この単純化は戦略的なものである。ボソンを消しフェルミオンを挿入する操作、あるいはその逆を行う操作を、エネルギーを変えずに理解したい。しかしその操作を構成する前に、同じ系を演算子の言語からラグランジアンの言語へ翻訳する。

\section{ボソン振動子}

ボソンに対して、実座標\(\phi(t)\)を導入する。空間変化がないので、自由スカラー場は通常の調和振動子になっている。黒板に最初に書かれた表式は
\[
\frac{1}{2}\dot\phi^{,2}+\frac{1}{2}m^2\phi^2.
\]
であった。この表式は運動エネルギーとポテンシャルエネルギーの和であり、ラグランジアンではない。符号は講義中に直ちに訂正される：
\begin{equation}
L_B=\frac{1}{2}\dot\phi^{,2}-\frac{1}{2}m^2\phi^2.
\label{eq:l7-boson-lagrangian}
\end{equation}

\begin{figure}[H]
\centering
\includegraphics[width=0.62\textwidth]{lecture_07_figure_02.png}
\lectureframe{00:03:25.431}
\caption{ボソンラグランジアンの上にある一モードハミルトニアン。このフレームは、\(T-V\)への口頭訂正の前の初期のプラス符号を保持している。}
\end{figure}

\begin{classroomqa}
\audiencequestion{表示された振動子ラグランジアンの符号の何が問題なのか？\lecturetimestamp{00:04:00}}
\lecturerresponse{エネルギーは\(T+V\)であるが、ラグランジアンは\(T-V\)である。したがってポテンシャル項はマイナス符号をもたなければならない。この訂正により、方程式は指数関数的振る舞いから意図された調和振動へと変わる。}
\end{classroomqa}

実際、オイラー--ラグランジュ方程式は
\begin{align}
\frac{d}{dt}\frac{\partial L_B}{\partial\dot\phi}
&=\frac{\partial L_B}{\partial\phi},\\
\ddot\phi&=-m^2\phi,
\end{align}
すなわち
\begin{equation}
\ddot\phi+m^2\phi=0.
\label{eq:l7-boson-eom}
\end{equation}
を与える。その振動数は\(m^2\)ではなく\(m\)である。この振動子を量子化すると、式~\eqref{eq:l7-one-mode-hamiltonian}の\(a^\dagger\)と\(a\)で記述されるのと同じ状態の梯子が得られる。

\section{なぜフェルミオン作用は1階なのか}

フェルミオンラグランジアンは時間微分を1つしか含まないため、より単純に見える。その理由は、意図的に素朴な試みの中にすでに見えている。通常の可換変数に対して、
\begin{equation}
L_{\mathrm{trial}}=\phi\dot\phi
=\frac{1}{2}\frac{d}{dt}\phi^2.
\end{equation}
これは全微分である。そのオイラー--ラグランジュ方程式は\(\dot\phi=\dot\phi\)と言うだけであり、したがってすべての微分可能な経路が許される。作用から全微分を取り除くことは端点のデータを変えるが、運動方程式の本体は変えない。

次に変数を\(\psi\)と改名し、それをグラスマン奇と宣言する。\(\psi\)と\(\dot\psi\)はともに反可換である。左変分の固定された選択のもとでは、奇因子の並べ替えがマイナス符号をもたらすので、同じ形式的な試みは
\begin{equation}
\frac{d}{dt}\frac{\partial L_{\mathrm{trial}}}{\partial\dot\psi}
=\dot\psi,
\qquad
\frac{\partial L_{\mathrm{trial}}}{\partial\psi}
=-\dot\psi.
\end{equation}
を与える。したがって方程式は
\begin{equation}
\dot\psi=-\dot\psi
\quad\Longrightarrow\quad
\dot\psi=0.
\label{eq:l7-constant-fermion}
\end{equation}
である。可換座標に対して恒真式であった表式が、反可換なものに対して力学的内容を獲得した。振動子にとっては内容が多すぎる。なぜなら定数しか許さないからであるが、これはまさに我々が必要とする手がかりである。

\begin{classroomqa}
\audiencequestion{解は単に\(\psi\)が定数であるということか？\lecturetimestamp{00:13:57}}
\lecturerresponse{はい。式~\eqref{eq:l7-constant-fermion}は真の1階微分方程式であり、その解は任意のグラスマン値の定数である。質量項があれば、その静的解は振動する解に変わる。}
\end{classroomqa}

明白な試みである\(m\psi^2\)は、\(\psi^2=0\)であるために失敗する。我々には二つの実奇変数\(\psi_1,\psi_2\)、あるいは等価的に一つの複素変数とその共役が必要である。講義では黒板で二実変数形式を丹念に進め、符号と因子2を繰り返し確認する。慣習的に規格化された複素形式は
\begin{equation}
L_F=\frac{i}{2}
\left(\bar\psi\dot\psi-\dot{\bar\psi}\psi\right)
-m\bar\psi\psi.
\label{eq:l7-fermion-lagrangian}
\end{equation}
であり、全微分を除いてこれは\(i\bar\psi\dot\psi-m\bar\psi\psi\)である。\(\bar\psi\)を変分すると
\begin{equation}
i\dot\psi=m\psi,
\qquad
\psi(t)=\psi(0)e^{-imt}.
\label{eq:l7-fermion-eom}
\end{equation}
が得られる。共役な方程式は逆の位相を持つ。したがってフェルミオン的モードは、その方程式が二階ではなく一階であるにもかかわらず、ボソン的モードとちょうど同じくエネルギー\(m\)で振動する。

\editorialnote{式~\eqref{eq:l7-fermion-lagrangian}は、探索的な二実変数計算に対する一つの標準的な順序と規格化を定める。講義自体は因子\(i\)を挿入し、因子2を調整し、一つの実成分の符号を未確定のまま残す。これらの変更は中心的な結論を変えない。すなわち、フェルミオンの運動は一階であり、質量は線形に現れ、複素フェルミオンは非零の双線形量を与える。}

これは、よく知られた相対論的対比のゼロ次元的な影である。Klein--Gordon方程式は時間について二階であり、Dirac方程式は一階である。空間を抑圧すると、ボソンは一つの調和振動子であり、フェルミオンは一つの複素Grassmann振動子である。

\begin{classroomqa}
\audiencequestion{\(\psi\)、\(\dot\psi\)、およびそれらの共役を書く順序は重要か。\lecturetimestamp{00:18:58}}
\lecturerresponse{中間の式は、奇因子が反可換であるために順序に依存する。物理的理論は、左または右の変分、共役、積分について一貫して単一の規約を用いる限り、任意の記述上の好みに依存しない。奇数の対を逆にすると符号が変わり、それは他のすべての場所でも一貫して変更されなければならない。}
\end{classroomqa}

式~\eqref{eq:l7-boson-lagrangian}と\eqref{eq:l7-fermion-lagrangian}を組み合わせると、式~\eqref{eq:l7-one-mode-hamiltonian}で表される同じ対になった組のラグランジアン記述が得られる。これで、それらの間の対称性を同定する準備が整った。

\section{超電荷とその代数}

フェルミオンを消去しボソンを挿入する演算子は
\begin{equation}
Q^\dagger=\sqrt{2m}\,a^\dagger c,
\end{equation}
であり、その共役はボソンを消去しフェルミオンを挿入する。
\begin{equation}
Q=\sqrt{2m}\,a c^\dagger.
\end{equation}
規格化は、代数が標準的な形をとるように選ばれている。\(cc^\dagger=1-c^\dagger c\)、\(aa^\dagger=a^\dagger a+1\)、およびボソン演算子とフェルミオン演算子が互いに可換であるという事実を用いると、
\begin{align}
\{Q,Q^\dagger\}
&=2m\left(ac^\dagger a^\dagger c+a^\dagger c a c^\dagger\right)\notag\\
&=2m\left[(a^\dagger a+1)c^\dagger c
+a^\dagger a(1-c^\dagger c)\right]\notag\\
&=2m\left(a^\dagger a+c^\dagger c\right)=2H.
\label{eq:l7-superalgebra}
\end{align}
ボソン演算子はフェルミオン演算子を通り抜けて移動させることができる。フェルミオン占有を制限するのと同じ冪零性により、
\begin{equation}
Q^2=(Q^\dagger)^2=0.
\end{equation}
が得られる。最後に、
\begin{equation}
[H,Q]=[H,Q^\dagger]=0.
\label{eq:l7-conserved-supercharges}
\end{equation}
各超電荷はボソン数とフェルミオン数を逆方向に変化させるので、振動数が等しいとき、二つのエネルギーの変化は相殺する。

\begin{classroomqa}
\audiencequestion{\([H,Q]\)とは何か。答えは\(Q\)に比例するのか。\lecturetimestamp{00:25:10}}
\lecturerresponse{いいえ。答えはゼロである。単独の生成演算子または消去演算子は、\(H\)との交換子においてそのエネルギーを獲得するであろうが、\(Q\)におけるボソン的寄与とフェルミオン的寄与は相殺する。その相殺こそがまさに等エネルギー条件である。}
\end{classroomqa}

\begin{classroomqa}
\audiencequestion{その交換子はすでに黒板に書かれていなかったか。\lecturetimestamp{00:25:45}}
\lecturerresponse{書かれていた。再発見は有用である。なぜなら、それは関係を単に写すのではなく解釈することを我々に強いるからである。すなわち、ハミルトニアンと可換であることは、その変換が対称性であり、その生成子が保存されることを意味する。}
\end{classroomqa}

\emph{電荷}という形容詞は、この保存則を反映している。しかしながら、\(Q\)を状態に付随する通常の同時測定可能な数として扱うべきではない。それはGrassmann奇であり、冪零であり、ボソン状態をフェルミオン状態へ写す。有用な言明は、それが縮退した超多重項内で作用する保存された生成子であるということである。
\editorialnote{この表現は、状態をエネルギーや電荷でラベル付けするという教室での類推を洗練する。個々の非エルミートな奇超電荷は、数値的な量子ラベルを持つ通常の観測量ではない。それらの保存された作用と完全な超代数の表現が、正しい分類を与える。}

\section{微分としての生成子}

代数は簡潔であるが、それだけでは相互作用する理論をどのように構築するかを教えてくれない。次の段階は幾何学的である。よく知られた対称性の生成子はすでに微分演算子として作用している。講義の規約では、
\begin{equation}
H=i\frac{\partial}{\partial t},
\qquad
P=-i\frac{\partial}{\partial x}.
\end{equation}

\begin{figure}[H]
\centering
\includegraphics[width=0.50\textwidth]{lecture_07_figure_03.png}
\lectureframe{00:28:12.720}
\caption{時間並進の生成子として表されたハミルトニアン。}
\end{figure}

角運動量も同様に、座標と微分の積である。超電荷が奇生成子であるなら、それらが奇座標に関する微分で表現できるかどうかを問うのは自然である。そこで時間軸を超空間
\begin{equation}
(t,\theta,\bar\theta),
\qquad
\theta^2=\bar\theta^2=0,
\qquad
\theta\bar\theta=-\bar\theta\theta.
\end{equation}
へと拡大する。この空間上の場は、通常の時間と、共役な一対のグラスマン座標に依存する。

この講義は、符号がいかに厄介であるかを意図的に明らかにする。複素共役から生じる符号、奇パラメータを移動させることから生じる符号、そして左微分ではなく右微分を選ぶことから生じる符号は、それぞれ異なる操作である。互いに相容れない黒板上の段階を保存するのではなく、ここで一つの規約を固定する。すべての奇微分は左微分であり、
\begin{equation}
q=\partial_\theta+i\bar\theta\partial_t,
\qquad
\bar q=\partial_{\bar\theta}+i\theta\partial_t.
\label{eq:l7-differential-supercharges}
\end{equation}
とする。すると
\begin{equation}
q^2=\bar q^{,2}=0,
\qquad
\{q,\bar q\}=2i\partial_t=2H.
\label{eq:l7-differential-algebra}
\end{equation}
となる。因子2は二つの交差項、\(\{\partial_\theta,\theta\}=1\) と \(\{\partial_{\bar\theta},\bar\theta\}=1\) から生じる。もし生成子が \(\partial_\theta\) と \(\partial_{\bar\theta}\) のみを含むなら、それらの反交換子は消えてしまい、ハミルトニアンを再現できない。

\begin{classroomqa}
\audiencequestion{\(q\) と \(\bar q\) における二つの時間微分項は同じ符号を持つべきか。\lecturetimestamp{00:31:01}}
\lecturerresponse{混合反交換子が \(2H\) となるように、一貫して選ばれなければならない。全体としてプラスの規約でも、それと一貫して共役をとった代替案でもよい。相対符号を決めるのは、記号 \(i\) についての視覚的直感ではなく、計算である。}
\end{classroomqa}

教室では繰り返し符号が確認され、その難しさは冗談になる。通勤中に符号が変わったように見える、というのである。この冗談は的を射ている。物理的な規約が何も変わっていなくても、グラスマン座標パラメータを座標を通して反交換させることで符号が変わることがある。式~\eqref{eq:l7-differential-supercharges} が、以下すべてに対する固定された基準である。

\subsection{座標のずれとしての超対称性}

\(\xi,\bar\xi\) を無限小の奇パラメータとする。超場に
\begin{equation}
\delta_\xi=\xi q+\bar\xi\bar q
\end{equation}
を作用させることは、上記の規約では、無限小座標変換
\begin{align}
\theta&\longrightarrow\theta+\xi,\\
\bar\theta&\longrightarrow\bar\theta+\bar\xi,\\
t&\longrightarrow t+i\left(\xi\bar\theta+\bar\xi\theta\right).
\label{eq:l7-coordinate-shift}
\end{align}
と等価である。したがってスカラー超場の変化は
\begin{equation}
\delta_\xi\Phi
=\left(\xi\partial_\theta+\bar\xi\partial_{\bar\theta}
+i\xi\bar\theta\partial_t+i\bar\xi\theta\partial_t\right)\Phi
=\left(\xi q+\bar\xi\bar q\right)\Phi.
\end{equation}
となる。これが超対称性の幾何学的な内容である。それは通常の時間の並進と結びついた、奇方向への並進である。二つの奇シフトは、式~\eqref{eq:l7-differential-algebra} のために時間並進に閉じる。

\begin{classroomqa}
\audiencequestion{論争の的となったマイナス符号は、時間を遡る運動を表しうるのか。\lecturetimestamp{00:58:13}}
\lecturerresponse{異なる符号規約は、表示された無限小時間シフトを逆転させうるが、別個の物理的な時間逆行過程が導入されているわけではない。重要なのは、座標変換と微分生成子が同じ代数に従うことである。}
\end{classroomqa}

\begin{classroomqa}
\audiencequestion{この拡張された時間には物理的実現があるのか、それとも単なる形式的な帳簿付けなのか。\lecturetimestamp{00:58:23}}
\lecturerresponse{有用な意味では帳簿付けである。すなわち、代数を幾何学にまとめ上げる。しかしそのまとめ上げは、相互作用が導入されると物理的に価値を持つようになる。超空間は、テンソル計算がローレンツスカラーの構成に役立つのとちょうど同じように、各段階でボソン--フェルミオン対称性を保つラグランジアンを構成するための系統的な計算手法を与える。}
\end{classroomqa}

\section{超場と補助成分}

奇座標が二つしかないため、偶超場の展開は有限で終わる：
\begin{equation}
\Phi(t,\theta,\bar\theta)
=\phi(t)+\bar\theta\psi(t)+\bar\psi(t)\theta
+D(t)\bar\theta\theta.
\label{eq:l7-superfield}
\end{equation}
偶単項式の係数 \(\phi\) と \(D\) はボソン的であり、奇単項式の係数 \(\psi\) と \(\bar\psi\) はフェルミオン的である。\(\theta^2\)、\(\bar\theta^2\)、またはそれ以上の項は残らない。時間依存性は展開しないままにしておく。\(t\) における通常のテイラー級数は終わらず、ここでは役に立たないからである。

\begin{classroomqa}
\audiencequestion{最上位の項における成分 \(D(t)\) とは何か。\lecturetimestamp{00:41:46}}
\lecturerresponse{この段階では、それは有限展開における最後のボソン的係数にすぎない。成分作用において、それは通常、運動項を持たず、代数方程式によって消去されるので、補助場と呼ばれる。講義では冗談めかして ``DはダミーのD'' と示唆しているが、これは歴史的な定義ではなく、その非力学的な役割をとらえたものである。}
\end{classroomqa}

超場は一つの対象であり、\(\phi,\psi,\bar\psi,D\) はその成分である。超対称性変換はそれらを混合する。したがって、ボソンとフェルミオンは、ベクトルの \(x\) 成分と \(y\) 成分が回転によって関係づけられるのとほぼ同じ構造的な意味で関係づけられる。ただし、グラスマン方向は通常の軸として思い描くことはできない。

\section{超作用と最上位成分}

運動方程式を得るには作用が必要である。通常の場の理論ではスカラーラグランジアンを構成し、それを時空にわたって積分する。ここでは対応する表式は
\begin{equation}
S=\int dt\,d\theta\,d\bar\theta\;\Lambda
\bigl(\Phi,\text{ superspace derivatives of }\Phi\bigr),
\label{eq:l7-superaction}
\end{equation}
であり、ここで\(\Lambda\)は偶スーパーフィールドである。ベレジン積分は代数的である。我々は
\begin{equation}
\int d\theta\,d\bar\theta\;\bar\theta\theta=1.
\end{equation}
を選ぶ。したがって、もし
\begin{equation}
\Lambda=A+\bar\theta B+\bar B\theta+C\bar\theta\theta,
\end{equation}
ならば
\begin{equation}
\int d\theta\,d\bar\theta\;\Lambda=C.
\label{eq:l7-top-component}
\end{equation}
である。奇積分は最上位成分を抽出し、時間にわたって積分されるべき通常のラグランジアンを残す。

なぜ結果は不変なのか。超対称性の変分は積分変数を平行移動する。最上位成分の変分は時間に関する通常の全微分であり、一方、奇の全微分はベレジン積分の下で消える。適切な端点条件の下で、作用は変わらない。これは時空にわたって積分する前にローレンツスカラーを構成することの超空間版である。

\section{スーパーフィールドの積}

古いスーパーフィールドから新しいスーパーフィールドを作る最も簡単な方法は、それをそれ自身と掛けることである。式~\eqref{eq:l7-superfield}から出発すると、\(\theta^2\)または\(\bar\theta^2\)を含むすべての項は消える。順序\(\bar\theta\theta\)を保ち、バー付き成分場を左に書くと
\begin{align}
\Phi^2
={}&\phi^2
+2\phi\bar\theta\psi
+2\phi\bar\psi\theta\notag\\
&+\left(2\phi D-2\bar\psi\psi\right)\bar\theta\theta.
\label{eq:l7-superfield-square}
\end{align}
となる。

\begin{figure}[H]
\centering
\includegraphics[width=0.82\textwidth]{lecture_07_figure_05.png}
\lectureframe{01:07:57.844}
\caption{成分ごとのスーパーフィールドの積とその\(\bar\theta\theta\)項の集まり。}
\end{figure}

二つの項\(\phi D\)と\(D\phi\)は、両方の場が偶であるため加算される。二つのフェルミオン交差項も、それぞれを並べ替えた後に加算されるが、それらの共通の符号は負である：
\begin{equation}
(\bar\theta\psi)(\bar\psi\theta)
+(\bar\psi\theta)(\bar\theta\psi)
=-2\bar\psi\psi\,\bar\theta\theta.
\end{equation}
この二乗を奇座標にわたって積分すると
\begin{equation}
\int d\theta\,d\bar\theta\;\Phi^2
=2\phi D-2\bar\psi\psi.
\end{equation}
が得られる。この項は構成により超対称であるが、まだ運動項ラグランジアンではない。時間微分を含まない。

次に板書は\(\Phi^3\)、\(\Phi^4\)、\(\Phi^n\)を演習として提案する。簡潔な規則がある。任意の通常の関数\(W\)に対して、
\begin{equation}
\left.W(\Phi)\right|_{\bar\theta\theta}
=W'(\phi)D-W''(\phi)\bar\psi\psi.
\label{eq:l7-function-top-component}
\end{equation}
である。\(W(\phi)=\phi^2\)に対しては、これは式~\eqref{eq:l7-superfield-square}を再現し、\(W(\phi)=\phi^n\)に対しては、
\begin{equation}
\left.\Phi^n\right|_{\bar\theta\theta}
=n\phi^{n-1}D-n(n-1)\phi^{n-2}\bar\psi\psi.
\end{equation}
を与える。
\editorialnote{講義は\(n\)乗を演習として割り当て、期待される項の種類を述べる。式~\eqref{eq:l7-function-top-component}は固定された順序の規約における完成した規則を与える。}

この時点で講師は会場が疲れ切っていないかと尋ねる。聴衆はいいえと答え、彼は自分が疲れ切っていると答える。このやり取りは単なる滑稽な息抜き以上のものである。スーパーフィールドの冪は展開が終了するため容易であるが、微分と拘束条件こそが符号の帳尻合わせが真に厳しくなる場所である。

\section{どの微分が超対称性を尊重するか？}

次の板書の議論を追う前に一つ訂正すべき点がある。通常の時間微分は無害である：
\begin{equation}
[\partial_t,q]=[\partial_t,\bar q]=0.
\end{equation}
したがって\(\partial_t\Phi\)は、\(\Phi\)がそうであるときはいつでもスーパーフィールドとして変換する。微妙な対象は裸の奇微分\(\partial_\theta\Phi\)と\(\partial_{\bar\theta}\Phi\)である。なぜならそれらは両方の超電荷と反可換でないからである。拘束条件および明白に超対称な成分構成のために修正された奇微分が必要である。
\editorialnote{時間微分が別のスーパーフィールドを生成しないという発言は、ここで用いられている規約では正しくない。\(\partial_t\)は両方の超電荷と可換であるため、共変である。修正された奇微分のその後の構成は依然として必要であり、講義の実際の数学的要点である。}

\subsection{``魔法の場''の類推}

講義は通常の回転を用いて基準を動機づける。\(f(x,y)\)が
\begin{equation}
\partial_y f=0.
\label{eq:l7-magic-constraint}
\end{equation}
を満たすとする。そのような関数を``魔法の場''と呼ぶ。それは\(x\)に依存するが\(y\)には依存しない。魔法の場のクラスは回転によって保たれるか？
\begin{equation}
L_z=-i\left(x\partial_y-y\partial_x\right),
\end{equation}
に対して
\begin{equation}
[\partial_y,L_z]=i\partial_x\neq0.
\end{equation}
が得られる。したがって\(L_z f\)は、\(f\)が最初に式~\eqref{eq:l7-magic-constraint}を満たしていても、一般に\(y\)依存性を獲得する。拘束条件は回転不変ではない。

一般的な教訓は単純である。場が\(C\Phi=0\)を満たし、対称性が\(G\)によって生成されるならば、\(C\)と\(G\)の次数付き交換子が消えるとき、変換された場は拘束されたままである。二つの奇演算子に対して、次数付き交換子は反交換子である。

\begin{classroomqa}
\audiencequestion{場に対する拘束条件が対称性と整合するのはいつか？\lecturetimestamp{01:17:56}}
\lecturerresponse{拘束演算子が対称性の生成子と次数付き可換であるときである。そのとき\(C(G\Phi)\)は\(G\)を通り過ぎて移動でき、\(G(C\Phi)=0\)に帰着する。回転の例は\(\partial_y\)が回転生成子と可換でないため失敗する。}
\end{classroomqa}

\subsection{共変な奇微分}

超電荷が式~\eqref{eq:l7-differential-supercharges}で固定されたので、
\begin{equation}
\mathcal D=\partial_\theta-i\bar\theta\partial_t,
\qquad
\bar{\mathcal D}=\partial_{\bar\theta}-i\theta\partial_t.
\label{eq:l7-covariant-derivatives}
\end{equation}
と定義する。唯一の変更は時間微分項の符号である。直接計算すると
\begin{align}
\{\mathcal D,q\}
&=\{\mathcal D,\bar q\}
=\{\bar{\mathcal D},q\}
=\{\bar{\mathcal D},\bar q\}=0,\\
\{\mathcal D,\bar{\mathcal D}\}&=-2i\partial_t.
\label{eq:l7-covariant-algebra}
\end{align}
が得られる。潜在的に非零となりうる2つのクロス項は相殺する。これは共変微分の奇数的類似物である。すなわち、変換則を保ちながら微分する。

講義では\(D\)を補助成分とこの微分の両方に用いている。我々は成分には\(D(t)\)を予約し、演算子にはカリグラフィックな\(\mathcal D\)を用いる。

\section{カイラル拘束条件}

ここで、超対称性のもとで生き残る非自明な拘束条件を課すことができる：
\begin{equation}
\bar{\mathcal D}\Phi=0.
\label{eq:l7-chiral-constraint}
\end{equation}
このような場はカイラルと呼ばれる。超対称性変換の前でカイラルであれば、
\begin{equation}
\bar{\mathcal D}(q\Phi)
=-q(\bar{\mathcal D}\Phi)=0,
\qquad
\bar{\mathcal D}(\bar q\Phi)
=-\bar q(\bar{\mathcal D}\Phi)=0.
\end{equation}
となる。したがって、変換された超場全体が拘束されたクラスの中に留まる。マジック場の条件とは異なり、カイラリティは対称性と両立する。

拘束条件を解くために、シフトされた時間座標を定義する：
\begin{equation}
\tau=t+i\bar\theta\theta.
\label{eq:l7-chiral-time}
\end{equation}
左微分を用いているので、
\begin{equation}
\partial_{\bar\theta}\tau=i\theta,
\end{equation}
であり、したがって
\begin{equation}
\bar{\mathcal D}F(\tau,\theta)
=\left(\partial_{\bar\theta}-i\theta\partial_t\right)F(\tau,\theta)
=i\theta\partial_\tau F-i\theta\partial_\tau F=0.
\end{equation}
となる。よって一般の局所カイラル超場は\(\bar\theta\)に\(\tau\)を通じてのみ依存する：
\begin{equation}
\Phi_{\mathrm{ch}}(\tau,\theta)
=\phi(\tau)+\bar\psi(\tau)\theta.
\label{eq:l7-chiral-field-tau}
\end{equation}

\begin{classroomqa}
\audiencequestion{\(\tau\)と\(\theta\)への依存だけがカイラル拘束条件を満たす唯一の形ですか？\lecturetimestamp{01:32:35}}
\lecturerresponse{局所的には、はい。微分方程式\(\bar{\mathcal D}\Phi=0\)は、独立な\(\bar\theta\)依存性がシフトされた座標\(\tau=t+i\bar\theta\theta\)に吸収されうることをまさに述べている。残る有限の展開は\(\tau\)の関数に\(\theta\)に比例する1つの項を加えたものである。}
\end{classroomqa}

元の時間座標で再展開すると、
\begin{align}
\phi(\tau)
&=\phi(t)+i\bar\theta\theta\dot\phi(t),\\
\bar\psi(\tau)\theta
&=\bar\psi(t)\theta,
\end{align}
となる。なぜなら第2行への補正は\(\theta^2\)を含むからである。したがって
\begin{equation}
\Phi_{\mathrm{ch}}(t,\theta,\bar\theta)
=\phi(t)+\bar\psi(t)\theta
+i\bar\theta\theta\dot\phi(t).
\label{eq:l7-chiral-field-components}
\end{equation}

\begin{classroomqa}
\audiencequestion{\(\phi(t+i\bar\theta\theta)\)のテイラー展開は2次の項を含むべきですか？\lecturetimestamp{01:33:51}}
\lecturerresponse{いいえ。最初の補正は\(i\bar\theta\theta\dot\phi\)であり、それより高次の補正はすべて\((\bar\theta\theta)^2=0\)を含む。聴衆はこれを即座に理解する。冪零性がすでに最初の2つの項を与え、級数を打ち切っているからである。}
\end{classroomqa}

これがカイラリティの経済性である。式~\eqref{eq:l7-superfield}の非拘束超場は2つのボソン成分スロットと2つのフェルミオン成分スロットをもつ。カイラル条件は、この最小の1次元例において独立な補助スロットを取り除き、最上成分を\(\dot\phi\)に結びつける。

\section{適切に停止する計算}

次に自然な操作は、カイラル場をその共役と組み合わせ、超空間上で積分することである。共役の展開を始めると、聴衆の一人が、\(\theta\bar\theta\)の順序は共役のもとで反転しなければならないのかと尋ねる。その時点で符号はもはや信頼できるようには見えない。はったりをかますよりも、計算は停止する。

\begin{classroomqa}
\audiencequestion{共役場の掛け算を後回しにしても、オチは得られますか？\lecturetimestamp{01:37:30}}
\lecturerresponse{はい。符号と共役を一貫して扱えば、カイラル場とその共役から作られた超空間表式は、望ましい通常のボソン的およびフェルミオン的運動項を生み出す。この講義にとって重要な結果は、疲れた計算から写し取った信頼できない最終行ではなく、構成原理である。}
\end{classroomqa}

停止点は見えるままにしておくべきである。この形式主義は容赦がない。1つの奇因子を置き間違えれば符号が変わり、同じ規約が共役、微分、掛け算、積分を通じて生き残らなければならない。その連鎖を失ったときの正しい対応は、先に進む前にそれを修復することである。標準的な規約では、おなじみの運動作用は共役積から作られるが、その詳細な成分還元は、未完成の黒板の行を再構成するのではなく、次のきれいな計算に属する。

\section{なぜ対称性がオチなのか}

超対称なラグランジアンは、ボソンとフェルミオンを並べるだけではない。それらが超代数の表現をなすように、質量と相互作用を拘束する。厳密に超対称な真空では、超多重項のメンバーは等しい質量をもつ。それらのループ寄与は相関し、ボソン的およびフェルミオン的紫外効果は、対称性によって強制されたまさにその組み合わせにおいて相殺する。

これが構成の仕組みが重要である理由である。ファインマン頂点を列挙し、図ごとに相殺を検証することもできるが、その場合相殺は偶然に見えてしまう。超対称性は、ローレンツ対称性があらゆる慣性系で散乱振幅が一貫して変換する理由を説明するのと同じ仕方で、それらを説明する。

\editorialnote{厳密な超対称性は、個々の自己エネルギー補正がすべて消えることを意味するものではない。それは繰り込みの間の関係を与え、特定の紫外発散への感受性を取り除く。非繰り込み性の主張は、量、次元、超対称性の量に依存する。講義で述べられた広い意味での相殺の主張は、この限定された形で保持される。}

\begin{classroomqa}
\audiencequestion{``chiral''はどう綴りますか。\lecturetimestamp{01:39:48}}
\lecturerresponse{C-H-I-R-A-L。この語は制約\(\bar{\mathcal D}\Phi=0\)を指す。三次元空間における物体の通常の幾何学的な左右性ではないが、用語自体は相対論的スピノルの左手系・右手系の構造に由来する。}
\end{classroomqa}

\begin{classroomqa}
\audiencequestion{ラグランジアンを経由せず、直接超対称ファインマン則に進むことはできますか。\lecturetimestamp{01:40:00}}
\lecturerresponse{規則を述べること自体は確かにできるが、その場合、結合定数間の関係やそれらが生み出す相殺は、説明のつかない偶然のように見える。より深い内容は一般化された時空対称性である。ラグランジアン、または等価な不変な定式化が、その理由を可視化し、規則を系統的に生成する。}
\end{classroomqa}

ローレンツ変換について何も言及していないのに、すべての摂動次数がたまたまローレンツ不変な振幅を与えたとしよう。十分な計算の後には、我々は説明を要求するだろう。答えは、すべての次数を一度に組織する時空対称性である。超空間の役割も同様である。\(\theta,\bar\theta\)におけるシフトはボソン的成分とフェルミオン的成分を混合し、それらの代数は時空の並進で閉じる。

\begin{classroomqa}
\audiencequestion{通常の空間座標も、グラスマン座標のように有限の冪級数に展開すべきですか。\lecturetimestamp{01:44:55}}
\lecturerresponse{いいえ。通常の座標における冪級数は終わらず、通常は自然な分解ではない。フーリエ変換の方が、場を波と運動量に分解するため、より有用である。\(\theta\)における展開が特別なのは、まさに冪零性によって有限になるからである。}
\end{classroomqa}

\begin{classroomqa}
\audiencequestion{もしLHCで新しい粒子が現れたら、それが単なる未知の粒子ではなく超対称パートナーであるとどうやって見分けられますか。\lecturetimestamp{01:45:42}}
\lecturerresponse{超対称性は、単に一つの余分な状態ではなく、相関した量子数と相互作用を予言する。スピンは半整数だけ異なり、ゲージ電荷は関係づけられ、多くの三次結合と四次結合は、一般的な理論よりも少ない独立パラメータで支配される。断面積、分岐比、およびスペクトルの残りの部分は、それらの関係を満たさなければならない。}
\end{classroomqa}

一つの見慣れない共鳴だけでは、超対称性を証明することにはならない。結合とスペクトルの関係のネットワークがそれを証明する。これが超空間構成の実験的な対応物である。一見独立な成分の相互作用が、一つの不変な対象から導かれるのである。

\begin{classroomqa}
\audiencequestion{観測されている粒子とその仮説上の超対称パートナーは等しい質量を持たない。これは理論をどのように変えますか。\lecturetimestamp{01:47:01}}
\lecturerresponse{それは理論を劇的に変える。超対称性は破れていなければならず、通常は自発的破れを通じて観測セクターに伝達される。破れは多重項の質量を分裂させるが、相互作用を制御し、スケールが十分低ければ紫外補正を制御するのに十分な構造を保持する。その機構を解明することは、特に重力を含めると、それ自体が相当な理論である。}
\end{classroomqa}

\section{単純な原理、困難な計算}

結びの評価は異例なほど率直である。出発点のアイデアは単純に聞こえる。ボソンとフェルミオンは対称性によって関係づけられる。しかしその実装は、たちまち添字、共役、グラスマン順序、符号の規律となる。対称性を構成した後でさえ、現実的な素粒子物理学にはそれを破ることが必要であり、重力と結合させると超重力へと向かい、さらに別の構造の層が加わる。

この困難は、そのアイデアが空虚である証拠ではない。自然は、単純な原理を計算しやすくする義務を負っていない。専門家は、符号の規則が第二の天性になるまで細部を繰り返し行うことでこの主題を学ぶ。不満は、簡潔な代数
\begin{equation}
\{Q,Q^\dagger\}=2H
\end{equation}
と、相互作用する理論でそれを実現するために必要な機械の量との間の隔たりから生じる。

\begin{classroomqa}
\audiencequestion{さらに別の抽象化の層が最終的に形式を単純化するのか、それとも複雑さが蓄積し続けるのか。\lecturetimestamp{01:49:33}}
\lecturerresponse{詳細な作業を取り除くような一般的な単純化は存在しない。基礎となる規則は短いかもしれないが、それらを信頼して使うには規約を繰り返し制御する必要がある。だからこそ、完全な扱いは通常、幅広い素粒子物理学のコースの一部ではなく、専門的なコース全体を占めるのである。}
\end{classroomqa}

\begin{classroomqa}
\audiencequestion{これらの変換は、フェルミオンとボソンが本当は何であるのか、あるいはなぜそれらの通常の振る舞いが異なるのかを教えてくれるのだろうか？\lecturetimestamp{01:51:42}}
\lecturerresponse{それらは、両者が一つの超場の成分であり、超電荷の下で互いの像であることを教えてくれる。\(Q\)を作用させるとボソン的成分がフェルミオン的成分に写され、次の変換は微分と補助成分を生み出す。これは回転がベクトルの成分を混合するのに似ているが、同じように文字通りの描像に翻訳することはできない。適切な言語はグラスマン幾何学そのものである。}
\end{classroomqa}

その答えは私たちを中心的な考えへと戻す。超対称性変換は超空間における微小な動きである。それは通常の回転が成分を混合するのと同じように\(\phi,\psi,\bar\psi,D\)を混合し、しかも時間の並進に閉じる。この描像は抽象的だが、恣意的ではない。超場のすべての項、共変微分のすべての符号、そして成分の相互作用間のすべての関係は、この代数を実現するために存在する。

\section{私たちが構築したもの}

私たちは対応する振動子ハミルトニアンから始め、その二つの異なるラグランジアン記述を回復した。ボソンは二階の振動子方程式に従い、フェルミオンは一階のグラスマン作用と奇成分の複素対を必要とする。超電荷は二つのモードを交換し、個々には二乗してゼロになり、反交換してハミルトニアンとなり、それと可換である。

それらの電荷を微分演算子として表すことで超空間に至った。有限の超場展開がボソン的、フェルミオン的、補助的成分を組み立て、ベレジン積分が超ラグランジアンの最上位成分を選び出し、乗法が不変な相互作用候補を生成した。対称性と両立する奇の微分の要求は、次に\(\mathcal D,\bar{\mathcal D}\)、カイラル拘束条件、そしてシフトされた座標\(\tau=t+i\bar\theta\theta\)へと導いた。

未完成の共役積の計算は、隠すべき欠陥ではなく教訓の一部である。超対称性は単純な物理的要求から始まるが、その計算はあらゆる順序の選択に敏感である。ひとたびそれらの選択が制御されれば、報酬は強力である。質量と結合は対称性によって組織され、量子補正は偶然でない関係を満たし、ボソンとフェルミオンは一つの幾何学的対象の成分となる。
